% Náhradní schéma stejnosměrného stroje s cizím buzením (motor).
% Kompilace: tectonic nahradni-schema.tex
%
% Tabulka spojení        součástka | svorka -> uzel
%   U   (napájení kotvy)   + -> A1,   - -> A2
%   R_a (odpor kotvy)      A1 -> m
%   U_i (indukované nap.)  + -> m,    - -> A2
%   U_f (napájení buzení)  + -> F1,   - -> F2
%   L_f (budicí vinutí)    F1 -> F2
\documentclass[border=8pt]{standalone}
\usepackage{fontspec}
\setmainfont{texgyreheros}[Extension=.otf, UprightFont=*-regular, BoldFont=*-bold, ItalicFont=*-italic]
\usepackage{amsmath}
\usepackage[european resistors,american inductors,american voltages,american currents]{circuitikz}
\ctikzset{bipoles/thickness=1.2}
\begin{document}
\begin{circuitikz}[line width=0.7pt]
% obvod kotvy
\coordinate (A1) at (0,3);  \coordinate (A2) at (0,0);
\coordinate (m)  at (3.6,3); \coordinate (mb) at (3.6,0);
% V se kreslí od svorky + ke svorce -
\draw (A1) to[V, l_={$U$}] (A2);
\draw (A1) to[R, l={$R_a$}, i>^={$I_a$}, *-] (m)
      to[V, l={$U_i$}] (mb) -- (A2);
\node[anchor=east, font=\small] at ($(A1)+(-0.12,0.28)$) {A1};
\node[anchor=east, font=\small] at ($(A2)+(-0.12,-0.28)$) {A2};
\draw (A2) node[circ]{};
% obvod buzení
\coordinate (F1) at (6.6,3); \coordinate (F2) at (6.6,0);
\draw (F1) to[V, l_={$U_f$}] (F2);
\draw (F1) -- (9.4,3) to[L, l={$L_f$}, i>^={$I_f$}] (9.4,0)
      -- (F2);
\draw (F1) node[circ]{} (F2) node[circ]{};
\node[anchor=east, font=\small] at ($(F1)+(-0.12,0.28)$) {F1};
\node[anchor=east, font=\small] at ($(F2)+(-0.12,-0.28)$) {F2};
\node[font=\small] at (1.8,-0.75) {obvod kotvy};
\node[font=\small] at (8.0,-0.75) {obvod buzení};
\end{circuitikz}
\end{document}
